\section{Phương pháp (Method)}
\label{sec:method}

Phần này mô tả toàn bộ pipeline đa tầng. Hình~\ref{fig:workflow} tổng hợp kiến trúc: từ một query khoa
học, hệ thống sinh candidate bằng các base retriever, (tùy chọn) chọn route, rồi dùng các tín hiệu bất
định rẻ tiền cùng một ngưỡng có bảo đảm conformal để quyết định \emph{khi nào} gọi cross-encoder reranking.

%================= TikZ WORKFLOW =================
\begin{figure}[t]
\centering
\resizebox{\textwidth}{!}{%
\begin{tikzpicture}[
  node distance=0.9cm and 1.2cm,
  font=\small,
  stage/.style={rectangle, rounded corners=3pt, draw=cstage!80, fill=cstage!12,
                text width=2.6cm, minimum height=1.0cm, align=center, thick},
  data/.style={rectangle, rounded corners=2pt, draw=cdata!80, fill=cdata!12,
               text width=2.4cm, minimum height=0.9cm, align=center},
  gate/.style={diamond, aspect=2, draw=cgate!85, fill=cgate!14,
               text width=2.0cm, align=center, inner sep=1pt, thick},
  arr/.style={-{Stealth[length=2.2mm]}, thick, draw=black!65},
  lbl/.style={font=\scriptsize\itshape, text=black!60}
]

% Input
\node[data] (q) {Query khoa học\\(claim)};

% Base retrievers
\node[stage, right=of q] (bm25) {BM25\\\scriptsize sparse lexical};
\node[stage, above=0.35cm of bm25] (dense) {Dense\\\scriptsize SciNCL};
\node[stage, below=0.35cm of bm25] (rrf) {Hybrid RRF\\\scriptsize $k{=}60$};

% Routing / candidate pool
\node[stage, right=1.6cm of bm25, text width=2.7cm] (cand) {Candidate pool\\Hybrid top-100};

% QPP signals
\node[stage, right=of cand, text width=2.8cm, fill=cstage!18] (qpp)
   {QPP signals\\\scriptsize Hybrid max score,\\ \scriptsize NQC, score-gap};

% Conformal gate
\node[gate, right=of qpp] (gate) {Conformal\\gate $\lambda(\alpha)$};

% Rerank vs skip
\node[stage, right=1.4cm of gate, text width=2.7cm, fill=cgate!12, draw=cgate!80] (rerank)
   {Cross-encoder rerank\\\scriptsize MiniLM-L-6-v2\\\scriptsize Hybrid top-20};
\node[data, below=1.1cm of rerank] (skip) {Giữ thứ tự\\Hybrid (skip)};

% Output
\node[data, right=1.2cm of rerank, text width=2.2cm] (out) {Ranked list\\cuối cùng};

% Arrows base
\draw[arr] (q) -- (bm25);
\draw[arr] (q.north) to[bend left=12] (dense.west);
\draw[arr] (q.south) to[bend right=12] (rrf.west);
\draw[arr] (dense) -- (cand);
\draw[arr] (bm25) -- (cand);
\draw[arr] (rrf) -- (cand);
\draw[arr] (cand) -- (qpp);
\draw[arr] (qpp) -- (gate);

% Gate decisions
\draw[arr, draw=cgate!85] (gate) -- node[lbl, above]{uncertain} (rerank);
\draw[arr, draw=cdata!75] (gate.south) |- node[lbl, below, pos=0.25]{confident} (skip);

\draw[arr] (rerank) -- (out);
\draw[arr] (skip.east) -| (out.south);

% Background groupings
\begin{scope}[on background layer]
  \node[draw=cstage!40, dashed, rounded corners, fit=(dense)(bm25)(rrf),
        inner sep=6pt, label={[lbl]above:Phase 1: Base retrieval}] {};
  \node[draw=cgate!40, dashed, rounded corners, fit=(qpp)(gate)(rerank)(skip),
        inner sep=8pt, label={[lbl]above:Phase 3: Uncertainty-aware selective reranking}] {};
\end{scope}

\end{tikzpicture}%
}
\caption{Workflow của hệ thống truy hồi đa tầng nhận biết bất định. Phase~1 sinh candidate bằng BM25,
Dense/SciNCL và Hybrid RRF. Phase~2 (query routing, không vẽ chi tiết để gọn) có thể chọn route cho từng
query. Phase~3 tính các tín hiệu \emph{QPP} rẻ tiền và dùng một \emph{conformal gate} với ngưỡng
$\lambda(\alpha)$ để quyết định: query \emph{uncertain} được đưa qua cross-encoder rerank trên Hybrid
top-20, query \emph{confident} giữ nguyên thứ tự Hybrid (skip). Cross-encoder --- tầng đắt nhất --- chỉ
chạy khi tín hiệu rẻ dự báo nó đáng giá.}
\label{fig:workflow}
\end{figure}

\subsection{Dataset và biểu diễn tài liệu}
\label{sec:method:data}
Các thí nghiệm dùng SciFact ở định dạng BEIR:
\begin{itemize}[leftmargin=1.6em,itemsep=1pt]
  \item Kích thước corpus: 5{,}183 documents.
  \item Test split: 300 queries. \quad Train split: 809 queries.
  \item Text sử dụng: chỉ \emph{title} và \emph{abstract}.
  \item Retrieval depth: top-100.
\end{itemize}
PDF toàn văn và passage chunking được cố ý loại trừ trong scope hiện tại. Điều này giữ retrieval unit
khớp với SciFact abstract và giữ thí nghiệm khả thi.

\subsection{Base retrievers}
\label{sec:method:base}
Hệ thống hiện thực ba base retrieval route:
\begin{description}[leftmargin=1.6em,style=nextline,itemsep=2pt]
  \item[BM25.] Một sparse lexical retriever dùng text title và abstract.
  \item[Dense.] Một SentenceTransformers dense retriever dùng \texttt{malteos/scincl} làm mô hình
  embedding khoa học chính.
  \item[Hybrid RRF.] Một hệ reciprocal-rank fusion trên các run BM25 và Dense/SciNCL, dùng $\texttt{rrf\_k}=60$.
\end{description}
Báo cáo dùng $\texttt{rrf\_k}=60$ như một setting RRF phổ biến cố định thay vì tune trên test set, nhằm giữ
baseline hybrid đơn giản và tránh thêm một nguồn test-set selection bias. Ablation các tham số RRF được để
lại cho future work.

\subsection{Oracle route labels}
\label{sec:method:oracle}
Với mỗi query, hệ thống tính per-query $\ndcg$ cho BM25, Dense và Hybrid. Oracle route label là route có
per-query $\ndcg$ tốt nhất. Các label này dùng để: (i) đo upper-bound của route selection; (ii) huấn luyện
và đánh giá router; (iii) khảo sát class imbalance và distribution shift. Phân bố oracle được trình bày ở
Bảng~\ref{tab:oracle-dist}; có một \emph{train/test distribution shift} rõ rệt: train thiên về Dense, còn
test thiên về BM25.

\begin{table}[h]
\centering
\caption{Phân bố oracle route label trên train và test split của SciFact.}
\label{tab:oracle-dist}
\begin{tabular}{lcc}
\toprule
Label & Train count & Test count \\
\midrule
BM25   & 195 & 226 \\
Dense  & 477 & 49  \\
Hybrid & 137 & 25  \\
\midrule
Tổng   & 809 & 300 \\
\bottomrule
\end{tabular}
\end{table}

\subsection{Các router baseline}
\label{sec:method:routers}
\begin{description}[leftmargin=1.6em,style=nextline,itemsep=2pt]
  \item[Random Router.] Lấy mẫu đều BM25, Dense, hoặc Hybrid.
  \item[Majority Router.] Luôn dự đoán nhãn oracle chiếm đa số (trên SciFact test là BM25).
  \item[Oracle Router.] Dùng trực tiếp oracle label; đây là upper bound.
  \item[Classical TF-IDF Logistic Regression Router.] Dùng TF-IDF features trên query text và balanced
  logistic regression, \emph{fit trên train split và predict trên test split} (disjoint, không leakage) ---
  một learned-router baseline hợp lệ song song với LLM router.
  \item[Small LLM QLoRA Router.] Dùng Qwen2.5-0.5B với QLoRA fine-tuning. Mô hình chỉ nhận query text và
  phải chọn đúng một nhãn: \texttt{bm25}, \texttt{dense}, hoặc \texttt{hybrid}.
\end{description}

\subsection{Label log-probability scoring}
\label{sec:method:logprob}
Router LLM đầu tiên dùng free-text generation và thường xuyên thất bại trong việc dự đoán BM25. Phiên bản
cải tiến chấm \emph{log probability} của ba nhãn route được phép và chọn nhãn có điểm cao nhất. Cách này
giảm lỗi định dạng output và khiến route confidence trở nên đo lường được.

\subsection{Calibration và confidence-gated fallback}
\label{sec:method:calib}
Calibrated router áp dụng: (i) một class bias cho mỗi nhãn route; (ii) temperature scaling trên label
score; (iii) một margin threshold giữa xác suất top-1 và top-2; (iv) fallback về Hybrid RRF khi label
margin dưới ngưỡng. Setting held-out tốt nhất trong grid compact hiện tại: temperature $=1.0$, margin
threshold $=0.3$, fallback $=$ Hybrid, biases $\text{BM25}=+0.75$, $\text{Dense}=+0.25$, $\text{Hybrid}=0.0$.

Vì không có file dev prediction sinh riêng từ Colab, thí nghiệm calibration tách file test prediction theo
query ID: 150 query cho calibration và 150 query rời rạc cho evaluation. Điều này tránh tune và evaluate
trên cùng query, nhưng vẫn chưa phải một setup train/dev/test hoàn toàn sạch.

\subsection{Retrieval-aware diagnostics}
\label{sec:method:diag}
Hệ thống cũng xuất các diagnostic nhẹ: BM25 / Dense / Hybrid top-score gap; BM25--Dense top-k overlap;
Hybrid--BM25 và Hybrid--Dense overlap; average RRF agreement. Các diagnostic này chưa phải feature mô hình
cuối cùng, nhưng giúp giải thích hành vi routing và nhận diện các tín hiệu bất định hữu ích (xem
Phần~\ref{sec:results:qpp}).

\subsection{Selective reranking}
\label{sec:method:rerank}
Phase~3 gồm: (i) một wrapper cross-encoder reranking với tự động phát hiện device CUDA/CPU; (ii) các tín
hiệu bất định; (iii) một trigger selective reranking.

Baseline \textbf{Always-Rerank} rerank Hybrid RRF top-20 candidate cho \emph{mọi} query bằng cross-encoder
\texttt{cross-encoder/ms-marco-MiniLM-L-6-v2}. Trên một GPU RTX~5070~Ti, nó chạy ở 27~ms/query, nhanh hơn
khoảng 14$\times$ so với CPU (khoảng 405~ms/query). Trigger \textbf{selective reranking} chỉ rerank các
query bị gắn cờ \emph{uncertain}. Vì reranking chỉ sắp xếp lại Hybrid top-20, ranking cuối của mỗi query
hoặc là thứ tự Hybrid gốc (skip) hoặc là thứ tự Always-Rerank (rerank); do đó một threshold sweep chỉ cần
trộn hai run đã tính sẵn ở mỗi operating point mà không cần thêm GPU inference.

\subsection{Query Performance Prediction làm tín hiệu trigger}
\label{sec:method:qpp}
Để tìm một trigger mạnh hơn nhưng vẫn label-free, chúng tôi tính các predictor QPP không giám sát trên các
run BM25, Dense và Hybrid: Weighted Information Gain (WIG), Normalized Query Commitment (NQC), độ lệch
chuẩn của top-k score, max score, và top-1/top-2 score gap, cộng với BM25--Dense top-10 disagreement. Mỗi
predictor được tương quan (Kendall $\tau$) với hai per-query target: base $\ndcg$, và \emph{gain-from-reranking}
(Always-Rerank $\ndcg$ trừ base $\ndcg$) --- chính là đại lượng mà một selective trigger cần dự báo.

\subsection{Conformal risk-controlled selective reranking}
\label{sec:method:conformal}
Thay vì chọn ngưỡng bằng cách tối đa hóa calibration $\ndcg$ (không có bảo đảm), chúng tôi áp dụng
\textbf{Conformal Risk Control} (CRC) để có bảo đảm finite-sample về chất lượng mất đi khi skip reranking.
Một query được rerank khi tín hiệu trigger của nó dưới ngưỡng; per-query loss là phần nDCG bị bỏ đi do
không rerank,
\begin{equation}
  \ell(q) = \max\!\big(0,\; \text{rerank\_ndcg}(q) - \text{base\_ndcg}(q)\big)
  \label{eq:loss}
\end{equation}
cho các query bị skip, và $0$ cho các query được rerank. Loss này đơn điệu theo ngưỡng, nên CRC chọn ngưỡng
coverage nhỏ nhất trên calibration split sao cho
\begin{equation}
  \frac{n \cdot \widehat{R} + B}{\,n + 1\,} \;\le\; \alpha,
  \qquad B = 1,
  \label{eq:crc}
\end{equation}
với $\widehat{R}$ là empirical risk và $B$ là loss bound, bảo đảm rằng kỳ vọng risk --- tức trung bình
phần nDCG thiếu hụt so với Always-Rerank --- tối đa là $\alpha$ trên held-out query.
